% !TeX program = lualatex
% !TeX encoding = utf8
% !TeX spellcheck = uk_UA
% !TeX root = ../EMConspect.tex

%=========================================================
\chapter{Випромінювання електромагнітних хвиль}\label{\currfilebase}
\graphicspath{{\currfilebase/Pictures}}
%=========================================================
\epigraph{Заряд, що спокійно рухається чи спочиває, надійно тримає своє поле при собі. Та щойно заряд смикнути --- прискорити, --- частина цього
	поля \enquote{відривається} від нього назавжди й вирушає в самостійну мандрівку простором зі швидкістю світла. Випромінювання --- це і є
	той момент, коли поле остаточно доводить свою незалежність від породжуючого його заряду.
}{Переосмислення ідей сучасної фізики}

%% --------------------------------------------------------
\section{Елементарна модель випромінювання}
%% --------------------------------------------------------

\subsection*{Випромінювання заряду: якісна картина}

Розглянемо електричне поле точкового заряду $q$. Якщо заряд перебуває в стані спокою, його електростатичне поле описується радіальними силовими
лініями, що виходять із заряду. Нехай у момент часу $t=0$ заряд під дією зовнішньої сили починає рухатися з прискоренням $a$, а через невеликий
час $\tau$ дія сили припиняється, після чого заряд рухається рівномірно зі швидкістю $v_q=a\tau$ (рис.~\ref{pic:VelocityGraph}); нехай при цьому
$v_q\ll c$, а $\tau$ --- дуже мале порівняно з часом спостереження $T$.

\begin{figure}[h!]\centering
	\localinput{velocity-graph.tikz}
	\caption{Швидкість заряду: короткий імпульс прискорення тривалістю $\tau$, після якого заряд
		рухається рівномірно\label{pic:VelocityGraph}}
\end{figure}

Оскільки збурення поля поширюється не миттєво, а зі швидкістю світла $c$, спостерігач, віддалений від заряду на різну відстань, дізнається про
прискорення заряду в різний момент часу. Нехай пройшов час $T\gg\tau$ (рис.~\ref{pic:ElementaryModel}):
\begin{itemize}
	\item \emphx{Зовнішній спостерігач} (на відстані, більшій за $cT$) іще нічого \enquote{не знає} про рух заряду взагалі --- сигнал про початок
	      руху до нього просто не встиг дійти, і він і далі бачить поле нерухомого заряду в його \emphx{початковому} положенні.
	\item \emphx{Внутрішній спостерігач} (на відстані, меншій за $c(T-\tau)$) уже \enquote{знає}, що прискорення закінчилось і заряд рухається
	      рівномірно, --- він бачить поле рівномірно рухомого заряду в його \emphx{поточному} положенні.
	\item У проміжній \emphx{тонкій сферичній оболонці} товщиною $c\tau$ --- \emphx{області випромінювання} --- силові лінії поля мусять
	      \enquote{зшивати} ці два різні поля (в різних точках простору поле відповідає різним, не пов'язаним безпосередньо положенням заряду),
	      утворюючи характерний \emphx{злам}.
\end{itemize}

\begin{figure}[h!]\centering
	\localinput{elementary-model.tikz}
	\caption{Заломлені (\enquote{злами}) силові лінії поля короткочасно прискореного заряду:
		межа між \enquote{старим} і \enquote{новим} полем\label{pic:ElementaryModel}}
\end{figure}

\begin{Attention}
	Саме цей злам силових ліній --- і є, по суті, електромагнітна хвиля, що \enquote{відірвалась} від заряду й розповсюджується самостійно.
	Оскільки лінії безперервні (не мають розривів, теорема Гаусса), вони не можуть просто обірватись на межі --- і саме необхідність \enquote{зшити}
	два різні поля без розриву силових ліній і породжує додаткову, \emphx{поперечну} складову поля, якої немає ні в чисто електростатичному, ні
	в чисто \enquote{рівномірно-рухомому} полі окремо.
\end{Attention}

\subsection*{Математичне виведення поля випромінювання}

Розглянемо детальніше злам силової лінії, що виходила з заряду під кутом $\theta$ до напряму його прискорення (рис.~\ref{pic:ElementaryDerivation}).
Нехай спостереження ведеться в момент $T=r/c$ (де $r$ --- відстань до точки спостереження), тобто поле, яке зараз \enquote{на підході}, --- це
поле, випромінене в момент $T-\tau$. Розкладемо поле $\Efield$ у зламі на \emphx{радіальну} (поздовжню) складову $E_\parallel$ і
\emphx{поперечну} складову $E_\perp$.

\begin{figure}[h!]\centering
	\localinput{elementary-model-derivation.tikz}
	\caption{Геометрія зламу силової лінії: розкладання поля на поздовжню $E_\parallel$ і
		поперечну $E_\perp$ складові\label{pic:ElementaryDerivation}}
\end{figure}

Із подібності трикутників (побудованих на бічних сторонах зламу) випливає:
\begin{equation}\label{eq:RatioEE}
	\frac{E_\perp}{E_\parallel} = \frac{v_qT\sin\theta}{c\tau}.
\end{equation}
Поздовжня складова --- це просто (майже незмінене) кулонівське поле заряду на відстані $r=cT$:
\begin{equation}\label{eq:EparDef}
	E_\parallel = \frac{q}{(cT)^2} \stackrel{r=cT}{=} \frac{q}{(cT)\,r}.
\end{equation}
Оскільки прискорення $a=v_q/\tau$, підставляючи в~\eqref{eq:RatioEE} і~\eqref{eq:EparDef}:
\begin{equation}\label{eq:EperpDerivation}
	E_\perp(t) = \frac{q}{(cT)\,r}\cdot\frac{v_qT\sin\theta}{c\tau} = \frac1{c^2}\frac{q}{r}\, a\,\sin\theta,
\end{equation}
де прискорення потрібно брати саме в той момент, коли воно й породило спостережуваний зараз злам, тобто в момент $t-r/c$:
\begin{equation}\label{eq:EperpResult}
	\tcbhighmath{
		E_\perp(t) = \frac1{c^2}\,\frac{q}{r}\, a\!\left(t-\frac{r}{c}\right)\sin\theta.
	}
\end{equation}

\begin{Attention}
	Формула~\eqref{eq:EperpResult} містить дві ключові риси \emphx{будь-якого} випромінювання, які ми надалі зустрічатимемо знову й знову:
	\begin{enumerate}
		\item Напруженість поля \emphx{хвилі} $E_\perp$ спадає як $1/r$, на відміну від електростатичного поля $E_\parallel\propto1/r^2$. Це
		      узгоджується із законом збереження енергії: енергія хвилі розподіляється по поверхні сфери ($\propto r^2$), а густина енергії
		      $\propto E^2$, тому $E^2\cdot r^2 = \const \Rightarrow E\propto1/r$.
		\item $E_\perp$ у момент часу $t$ визначається прискоренням заряду $a$ не в цей самий момент, а в \emphx{більш ранній} момент $t-r/c$
		      --- \emphx{час запізнення}, потрібний хвилі, щоб дістатися точки спостереження на відстані $r$.
	\end{enumerate}
\end{Attention}

%% --------------------------------------------------------
\section{Дипольне випромінювання}
%% --------------------------------------------------------

\subsection*{Диполь як елементарний випромінювач}

\begin{Attention}
	\emphx{Монопольного випромінювання не існує!} Ізольований заряд не може \enquote{пульсувати} (породжувати хвилю), бо це порушило б закон
	збереження заряду: сумарний заряд системи --- величина стала, а отже, і поле на великих відстанях (визначене насамперед повним зарядом) не
	може містити хвильової, змінної в часі складової від самого лише повного заряду. Найпростішою системою, здатною ефективно випромінювати, є
	\emphx{електричний диполь} зі змінним у часі дипольним моментом.
\end{Attention}

Розглянемо електронейтральну систему з двох зарядів $+q$ і $-q$, з яких додатний --- нескінченно важкий (і тому нерухомий), а від'ємний вільно
рухається (рис.~\ref{pic:DipoleTwoCharges}). Дипольний момент такої системи та його друга похідна за часом:
\begin{equation*}
	\vect p = -q\vect r', \qquad \ddot{\vect p} = -q\dot{\vect v},
\end{equation*}
де $\vect v=\dot{\vect r}'$ --- швидкість заряду, а $\dot{\vect v}$ --- його прискорення.

\begin{figure}[h!]\centering
	\localinput{dipole-two-charges.tikz}
	\caption{Найпростіша модель диполя: важкий нерухомий заряд $+q$ і легкий рухомий $-q$\label{pic:DipoleTwoCharges}}
\end{figure}

\subsection*{Про запізнюючі потенціали (коротко)}

Строге виведення полів диполя спирається на розв'язання рівнянь Максвелла через електромагнітні потенціали $\phi$, $\vect A$ у калібруванні
Лоренца, що приводить до неоднорідного хвильового рівняння Даламбера з розв'язками у вигляді \emphx{запізнюючих потенціалів}:
\begin{equation*}
	\vect A(\vect r, t)= \frac1c \iiint\limits_{V'} \frac{\vect j\!\left(\vect r', t - \dfrac{|\vect r-\vect r'|}{c}\right) dV'}{|\vect r-\vect
	r'|}, \qquad \phi(\vect r, t) = \iiint\limits_{V'} \frac{\rho\!\left(\vect r', t - \dfrac{|\vect r-\vect r'|}{c}\right) dV'}{|\vect r-\vect
	r'|}.
\end{equation*}

\begin{Attention}
	Ці потенціали називають \emphx{запізнюючими}, бо поле в точці $\vect r$ у момент $t$ визначається не миттєвим розподілом зарядів і струмів, а
	їхнім станом у \emphx{попередній} момент часу $t-|\vect r-\vect r'|/c$ --- з урахуванням скінченної швидкості поширення поля. Формально
	рівняння Даламбера допускають і другий розв'язок, зі знаком \enquote{$+$} в аргументі (\emphx{випереджувальні потенціали}), --- але він не має
	ясного фізичного сенсу (поле випереджало б причину, яка його породжує) і відкидається на основі \emphx{причинності}. Розкладаючи ці інтеграли
	в ряд для випадку $r\gg r'$ (система набагато менша за відстань до спостерігача) і нерелятивістських швидкостей зарядів ($v\ll c$), можна
	довести, що головний член розкладу визначається виключно \emphx{дипольним моментом} системи $\vect p(t-r/c)$ --- це й строго обґрунтовує, чому
	диполь є найпростішим (і найважливішим) елементарним випромінювачем; наступні члени розкладу відповідають магнітно-дипольному,
	електрично-квадрупольному й вищим типам випромінювання, зазвичай значно слабшим за дипольне.
\end{Attention}

\subsection*{Поля диполя: ближня й хвильова зони}

Нехай дипольний момент змінюється гармонічно, $\vect p = \vect p_0\cos\omega t$. Точне вираження полів (у сферичних координатах, з віссю $z$
уздовж $\vect p_0$) містить доданки різного порядку за $1/r$, і його зручно аналізувати окремо у двох граничних випадках.

\emphx{Ближня зона} --- відстань від осцилятора до точки спостереження мала порівняно з довжиною хвилі:
\begin{equation}\label{eq:NearZoneCondition}
	r \ll \frac\lambda{2\pi}.
\end{equation}
У цій зоні електричне поле в кожен момент часу $t$ практично збігається з полем \emphx{статичного} диполя, момент якого дорівнює миттєвому
значенню $p(t)$:
\begin{equation}\label{eq:NearZoneE}
	E_r = \frac{2\cos\theta\, p(t)}{r^3}, \qquad E_\theta = \frac{\sin\theta\, p(t)}{r^3},
\end{equation}
а магнітне поле, оскільки $\dot p(t) = Il$ (де $I$ --- еквівалентний струм, $l$ --- довжина еквівалентного елемента струму), збігається з
квазістатичним полем цього елемента струму за законом Біо--Савара--Лапласа:
\begin{equation}\label{eq:NearZoneB}
	\Bfield = \frac1{cr^3}\left[\dot{\vect p}(t)\times\vect r\right].
\end{equation}

\begin{Attention}
	У ближній зоні поля \emphx{не запізнюються} (запізнення $r/c \ll T = 2\pi/\omega$ --- мізерне порівняно з періодом коливань) і практично
	збігаються з полями відповідного статичного диполя й елемента струму --- випромінювання тут ще \enquote{не встигло сформуватися}.
\end{Attention}

\emphx{Хвильова зона} --- відстань, набагато більша за довжину хвилі:
\begin{equation}\label{eq:WaveZoneCondition}
	r \gg \frac\lambda{2\pi}.
\end{equation}
У цій зоні (рис.~\ref{pic:WaveZone}) залишаються лише члени, що спадають як $1/r$ (решта --- швидше, і на великих відстанях нехтовно малі):
\begin{equation}\label{eq:WaveZoneFields}
	\tcbhighmath{
		E_\theta = B_\varphi = \frac{\sin\theta}{c^2r}\,\ddot p\!\left(t-\frac{r}c\right), \qquad E_r=E_\varphi=B_r=B_\theta=0.
	}
\end{equation}

\begin{figure}[h!]\centering
	\includegraphics[width=0.45\linewidth]{wave_zone}
	\caption{Поля $E_\theta$ і $B_\varphi$ у хвильовій зоні диполя\label{pic:WaveZone}}
\end{figure}

\begin{Attention}
	У хвильовій зоні поля $\Efield$ і $\Bfield$ поперечні, чисельно рівні одне одному й спадають як $1/r$ --- точнісінько ті самі риси плоскої
	електромагнітної хвилі (гл.~\ref{EMWaves})! Крім того, поле залежить від полярного кута $\theta$: на осі диполя ($\theta=0,\pi$) поле
	дорівнює нулю, максимуму ж воно сягає в екваторіальній площині ($\theta=\pi/2$) --- \emphx{диполь не випромінює вздовж власної осі}. У кожній
	точці хвильової зони вектори $\Efield$, $\Bfield$ і $\vect r$ взаємно перпендикулярні й утворюють правовинтову трійку.
\end{Attention}

\begin{figure}[h!]
	\centering
	\begin{minipage}[t]{0.45\linewidth}\centering
		\includegraphics[width=\linewidth]{dipole_radiation}\\
		\footnotesize Ближня й хвильова зони: поле \enquote{відривається} від диполя
	\end{minipage}\hfill
	\begin{minipage}[t]{0.45\linewidth}\centering
		\includegraphics[width=\linewidth]{Fields}\\
		\footnotesize Миттєвий знімок силових ліній $\Efield$ (сині) і $\Bfield$ (червоні)
	\end{minipage}
	\caption{Поле осцилюючого диполя: у ближній зоні поле ніби \enquote{причеплене} до диполя, у
		хвильовій --- випромінюється\label{pic:DipoleRadiationSnapshot}}
\end{figure}

\subsection*{Діаграма спрямованості}

Кутову залежність інтенсивності випромінювання диполя ($\propto\sin^2\theta$) зручно зображати \emphx{полярною діаграмою спрямованості}
(рис.~\ref{pic:PolarPattern}): довжина відрізка від початку координат до кривої в напрямку $\theta$ дорівнює величині поля (чи інтенсивності) у
цьому напрямку.

\begin{figure}[h!]
	\centering
	\begin{minipage}[t]{0.45\linewidth}\centering
		\localinput{polar-pattern.tikz}
	\end{minipage}\hfill
	\begin{minipage}[t]{0.45\linewidth}\centering
		\includegraphics[width=\linewidth]{polar_diagramm}
	\end{minipage}
	\caption{Полярна діаграма спрямованості дипольного випромінювача: характерна форма
		\enquote{бублика}\label{pic:PolarPattern}}
\end{figure}

%% --------------------------------------------------------
\section{Потужність випромінювання осцилятора}
%% --------------------------------------------------------

Густина потоку електромагнітної енергії характеризується вектором Пойнтінга $\vect\Pi$ (гл.~\ref{MaxwellEqns}); потік енергії крізь поверхню
сфери радіуса $r$, що оточує осцилятор, і називають \emphx{потужністю випромінювання}. Елемент площі сфери в сферичних координатах $d\sigma =
r^2\sin\theta\, d\theta\, d\varphi$, тому, підставляючи поля хвильової зони~\eqref{eq:WaveZoneFields} у $\Pi = \dfrac{c}{4\pi}EB = \dfrac{c}{4\pi}E^2$:
\begin{equation}\label{eq:PowerInstant}
	\frac{dW}{dt} = \oiint\limits_\sigma \Pi\, d\sigma = \frac{c}{4\pi}\int\limits_0^{2\pi}\!\!\int\limits_0^\pi \left[\frac{\sin\theta}{c^2r}\,
	\ddot p\!\left(t-\frac{r}{c}\right)\right]^2 r^2\sin\theta\, d\theta\, d\varphi = \frac{\ddot p^2\!\left(t-\dfrac{r}{c}\right)}{4\pi
	c^3}\cdot2\pi\int\limits_0^\pi\sin^3\theta\, d\theta.
\end{equation}
Оскільки $\displaystyle\int_0^\pi\sin^3\theta\,d\theta = \frac43$, отримуємо
\begin{equation}\label{eq:PowerInstantResult}
	\tcbhighmath{
		\frac{dW}{dt} = \frac{2\ddot p^2\!\left(t-\dfrac{r}{c}\right)}{3c^3}.
	}
\end{equation}
Для гармонічного осцилятора $p=p_0\cos\omega t$, $\ddot p = -\omega^2p_0\cos\omega t$, і, усереднюючи за період ($\langle\cos^2\rangle=1/2$):
\begin{equation}\label{eq:PowerAveraged}
	\tcbhighmath{
		\left\langle\frac{dW}{dt}\right\rangle = \frac{\omega^4p_0^2}{3c^3}.
	}
\end{equation}

\begin{Attention}
	Середня потужність випромінювання осцилятора пропорційна квадрату амплітуди дипольного моменту $p_0$ і --- значно різкіше --- \emphx{четвертій
		степені частоти} $\omega$. Саме тому для радіозв'язку використовують відносно \emphx{короткі} хвилі (десятки метрів --- кілометри, тобто
	порівняно високі частоти за мірками радіодіапазону): довгі хвилі (наднизькі частоти) від повільно змінних струмів практично не
	випромінюються --- $\omega^4$ у знаменнику зростання частоти дає надзвичайно різкий виграш у потужності випромінювання. Реальні радіоантени
	часто мають розмір порядку половини довжини хвилі --- саме такі \emphx{піврильні (half-wave) диполі} (рис.~\ref{pic:HalfWaveDipole})
	найефективніше випромінюють і приймають електромагнітні хвилі.
\end{Attention}

\begin{figure}[h!]\centering
	\includegraphics[width=0.5\linewidth]{HalfWaveDipole}
	\caption{Реальна піврильна дипольна антена\label{pic:HalfWaveDipole}}
\end{figure}

\subsection*{Чому небо блакитне}

Тим же характером залежності випромінювання осцилятора від частоти пояснюється, зокрема, і блакитний колір неба (рис.~\ref{pic:BlueSky}).
Сонячне світло, що пронизує атмосферу, розсіюється молекулами повітря, які можна вважати елементарними осциляторами: під впливом світлової хвилі
вони здійснюють \emphx{вимушені} коливання. Оскільки період власних коливань цих молекул істотно відрізняється від періоду видимого світла
(резонанс відсутній), амплітуда вимушених коливань слабко залежить від частоти хвилі, тому інтенсивність розсіяного світла, тобто інтенсивність
випромінювання цих вимушених осциляторів, за формулою~\eqref{eq:PowerAveraged}, обернено пропорційна $\lambda^4$ (позаяк $\omega\propto1/\lambda$).
Тому короткохвильове (синє) світло розсіюється значно сильніше, ніж довгохвильове (червоне), --- і небо здається блакитним.

\begin{figure}[h!]\centering
	
\includegraphics[width=0.4\linewidth]{BlueSky}
	\caption{Розсіювання сонячного світла молекулами атмосфери --- причина блакитного кольору
		неба\label{pic:BlueSky}}
\end{figure}

%% --------------------------------------------------------
\section{Випромінювання рухомого заряду}
%% --------------------------------------------------------

Скориставшись моделлю \enquote{важкий $+q$ --- легкий рухомий $-q$} (рис.~\ref{pic:DipoleTwoCharges}) і формулою~\eqref{eq:PowerInstantResult},
підставимо $\ddot{\vect p} = -q\dot{\vect v}$:
\begin{equation}\label{eq:LarmorFormula}
	\tcbhighmath{
		\frac{dW}{dt} = \frac{2q^2}{3c^3}\dot{\vect v}^2.
	}
\end{equation}

\begin{Attention}
	Формула~\eqref{eq:LarmorFormula} (\emphx{формула Лармора}) --- один із найважливіших результатів класичної електродинаміки: \emphx{будь-який}
	заряд, що рухається з прискоренням, неминуче випромінює електромагнітну енергію, причому втрачена потужність залежить \emphx{лише} від
	заряду й миттєвого прискорення --- і зовсім не залежить від того, чому саме заряд прискорюється (електричне поле, магнітне поле, механічний
	поштовх). Заряд, що рухається \emphx{рівномірно} й прямолінійно ($\dot{\vect v}=0$), не випромінює взагалі --- випромінює лише
	\emphx{прискорений} заряд.
\end{Attention}

\begin{Attention}
	Формула Лармора має драматичний історичний наслідок: електрон, що обертається по коловій орбіті навколо ядра (як у класичній
	\enquote{планетарній} моделі атома Резерфорда), рухається з доцентровим прискоренням, а отже, за~\eqref{eq:LarmorFormula}, \emphx{неминуче
		випромінює} й безперервно втрачає енергію. Оцінка (аналогічна до задачі~3 нижче) показує, що класичний електрон мав би \emphx{впасти на ядро}
	за час порядку $10^{-11}$~с --- атоми в такій моделі просто не могли б бути стабільними! Саме ця суперечність між класичною електродинамікою
	й спостережуваною стабільністю атомів стала одним з головних поштовхів до створення квантової механіки: М.~Бор 1913~р. був змушений
	постулювати, що на деяких (\enquote{стаціонарних}) орбітах електрон \emphx{всупереч} класичній електродинаміці не випромінює --- цей
	постулат, який не мав жодного пояснення в межах класичної фізики, згодом отримав строге обґрунтування вже в повній квантовій теорії.
\end{Attention}

%% --------------------------------------------------------
\section{Задачі}
%% --------------------------------------------------------

\subsection*{Задача 1}

\emphx{Умова.} Електромагнітна хвиля, що випромінюється диполем, поширюється у вакуумі так, що у хвильовій зоні на промені, перпендикулярному до
осі диполя, на відстані $r$ від нього середнє значення густини потоку енергії дорівнює $\langle\Pi\rangle$. Знайти середню потужність
випромінювання диполя.

\emphx{Розв'язання.} Оскільки $\Pi\propto E^2\propto\sin^2\theta$ (формула~\eqref{eq:WaveZoneFields}), а на перпендикулярному промені
$\theta=\pi/2$, тобто $\langle\Pi(\theta)\rangle = \langle\Pi\rangle\sin^2\theta$. Інтегруючи по сфері (як і при виведенні~\eqref{eq:PowerInstant}):
\begin{equation}\label{eq:Problem1Solution}
	\left\langle\frac{dW}{dt}\right\rangle = \oiint\limits_\sigma \langle\Pi(\theta)\rangle\, d\sigma = \langle\Pi\rangle\, r^2\cdot2\pi\int\limits_0^\pi
	\sin^3\theta\, d\theta = \frac{8\pi r^2}{3}\langle\Pi\rangle.
\end{equation}

\subsection*{Задача 2}

\emphx{Умова.} Сталий за модулем електричний дипольний момент $\vect p$ обертають з кутовою швидкістю $\omega$ навколо осі, перпендикулярної до
осі диполя й що проходить через його середину. Знайти потужність випромінювання диполя.

\emphx{Розв'язання.} Розкладемо вектор диполя на дві взаємно перпендикулярні складові в площині обертання: $p_x=p\cos\omega t$,
$p_y=p\sin\omega t$. Тоді
\begin{equation*}
	\ddot p_x = -p\omega^2\cos\omega t, \qquad \ddot p_y = -p\omega^2\sin\omega t, \qquad \ddot{\vect p}^{\,2} = \ddot p_x^2+\ddot p_y^2 =
	p^2\omega^4\left(\cos^2\omega t+\sin^2\omega t\right) = p^2\omega^4 = \const.
\end{equation*}
За формулою~\eqref{eq:PowerInstantResult} (яка тут одразу постійна в часі, тож усереднення не потрібне):
\begin{equation}\label{eq:Problem2Solution}
	\tcbhighmath{
		\frac{dW}{dt} = \frac{2p^2\omega^4}{3c^3}
	}
	--- \text{рівно вдвічі більше, ніж у лінійного осцилятора~\eqref{eq:PowerAveraged} з тією самою амплітудою $p$.}
\end{equation}

\begin{Attention}
	Цей результат має простий фізичний зміст: обертовий диполь --- це суперпозиція \emphx{двох} лінійних осциляторів у взаємно перпендикулярних
	напрямках, зсунутих за фазою на $90^\circ$ (так само, як обертовий рух по колу --- суперпозиція двох перпендикулярних гармонічних коливань).
	Оскільки миттєва потужність кожного окремого осцилятора коливається в протифазі одна до одної (коли $\cos^2\omega t$ максимальний,
	$\sin^2\omega t$ --- мінімальний, і навпаки), їхня \emphx{сума} виявляється сталою в часі, а не лише в середньому, --- обертовий диполь
	випромінює цілком рівномірно, без пульсацій потужності.
\end{Attention}

\subsection*{Задача 3}

\emphx{Умова.} Знайти потужність випромінювання нерелятивістської частинки із зарядом $e$ і масою $m$, що рухається коловою орбітою радіуса $R$
у полі нерухомого точкового заряду $q$.

\emphx{Розв'язання.} Кулонівська сила забезпечує доцентрове прискорення: $\dfrac{eq}{R^2} = \dfrac{mv^2}{R}$, звідки доцентрове прискорення
\begin{equation*}
	\dot v = \frac{v^2}{R} = \frac{eq}{mR^2}.
\end{equation*}
Підставляючи у формулу Лармора~\eqref{eq:LarmorFormula}:
\begin{equation}\label{eq:Problem3Solution}
	\tcbhighmath{
		\frac{dW}{dt} = \frac{2e^2}{3c^3}\left(\frac{eq}{mR^2}\right)^2 = \frac{2e^4q^2}{3c^3m^2R^4}.
	}
\end{equation}

\begin{Attention}
	Саме ця формула (з $e=q=$ заряд електрона, $m=$ маса електрона, $R\sim10^{-10}$~м --- розмір атома) і дає ту саму драматичну оцінку часу
	падіння класичного електрона на ядро $\sim10^{-11}$~с, про яку йшлося вище: досить прирівняти втрату енергії $dW/dt$ до швидкості зменшення
	повної механічної енергії електрона на орбіті, що зменшується разом з $R$, і проінтегрувати від початкового радіуса орбіти до нуля.
\end{Attention}

%% --------------------------------------------------------
\section{Підсумки розділу}
%% --------------------------------------------------------

\begin{table}[h!]\small\centering
	\caption{Основні співвідношення теорії випромінювання}
	\begin{tabular}{ll}
		\toprule
		Поле хвильової зони диполя      & $\displaystyle E_\theta = B_\varphi = \frac{\sin\theta}{c^2r}\,\ddot p\!\left(t-\frac rc\right)$ \\
		Миттєва потужність випромінювання диполя & $\displaystyle \frac{dW}{dt} = \frac{2\ddot p^2}{3c^3}$                                 \\
		Середня потужність гармонічного осцилятора & $\displaystyle \left\langle\frac{dW}{dt}\right\rangle = \frac{\omega^4p_0^2}{3c^3}$     \\
		Формула Лармора (довільний заряд) & $\displaystyle \frac{dW}{dt} = \frac{2q^2}{3c^3}\dot{\vect v}^2$                                \\
		\bottomrule
	\end{tabular}
\end{table}

\begin{Attention}
	Розділ про випромінювання замикає весь курс електрики й магнетизму найглибшим можливим підсумком: \emphx{будь-який} прискорений заряд
	неминуче породжує хвилю, що назавжди йде від нього в навколишній простір, --- а електромагнітне поле остаточно постає перед нами не як
	допоміжна математична конструкція для опису сил між зарядами, а як цілком самостійна, фізично реальна субстанція, здатна існувати,
	поширюватись і переносити енергію \emphx{без} будь-яких зарядів поблизу. Саме на цьому явищі --- випромінюванні прискорених зарядів ---
	ґрунтується вся радіотехніка, вся астрофізика випромінювання зір і туманностей, і, зрештою, саме воно --- а не якийсь окремий, спеціально
	влаштований \enquote{носій світла} --- відповідає за існування й поширення видимого світла в природі.
\end{Attention}
